\documentclass[12pt,a4paper]{article}
\usepackage{fontspec}
\setmainfont{Times New Roman}[Path=C:/Windows/Fonts/,UprightFont=times.ttf,BoldFont=timesbd.ttf,ItalicFont=timesi.ttf,BoldItalicFont=timesbi.ttf]
\usepackage{amsmath,amssymb,amsfonts}
\usepackage[left=1.5cm,right=1.5cm,top=1.5cm,bottom=1.7cm,headheight=17pt]{geometry}
\usepackage{enumitem,array,fancyhdr,booktabs}

\pagestyle{fancy}
\fancyhf{}
\lhead{\small\textbf{ĐÁP ÁN ĐỀ CƯƠNG CHƯƠNG V (ĐỀ 2) -- PHƯƠNG PHÁP TỌA ĐỘ TRONG KHÔNG GIAN}}
\rhead{\small\textit{Biên soạn: Dương Gia Bảo}}
\rfoot{\small Trang \thepage}
\renewcommand{\headrulewidth}{0.4pt}
\fancypagestyle{firstpage}{\fancyhead{}\renewcommand{\headrulewidth}{0pt}}

\setlength{\parindent}{0pt}
\setlength{\parskip}{4pt}
\setlength{\emergencystretch}{1em}

\setlist[enumerate,1]{label=\textbf{Câu \arabic*:},leftmargin=*,itemsep=7pt,topsep=4pt}

\begin{document}
\thispagestyle{firstpage}

\noindent
\begin{minipage}[t]{0.36\textwidth}
    \centering
    {\fontsize{12}{14}\selectfont\textbf{ĐỀ CƯƠNG ÔN TẬP}}\\[4pt]
    \textit{Biên soạn: Dương Gia Bảo}\\[6pt]
    \setlength{\fboxrule}{0.8pt}
    \setlength{\fboxsep}{5pt}
    \fbox{\begin{minipage}{0.88\linewidth}
        \centering
        {\fontsize{12}{14}\selectfont\textbf{ĐÁP ÁN CHƯƠNG V -- ĐỀ 2}}
    \end{minipage}}
\end{minipage}\hfill
\begin{minipage}[t]{0.62\textwidth}
    \centering
    {\fontsize{13}{15}\selectfont\textbf{ĐÁP ÁN \& HƯỚNG DẪN GIẢI CHI TIẾT}}\\[3pt]
    {\fontsize{12}{14}\selectfont\textbf{CHƯƠNG V: PHƯƠNG PHÁP TỌA ĐỘ TRONG KHÔNG GIAN (ĐỀ 2)}}\\[3pt]
    {\fontsize{11}{13}\selectfont\textit{(GDPT 2018 -- 20 Trắc nghiệm, 5 Đúng/Sai, 25 Trả lời ngắn)}}
\end{minipage}

\par\vspace{4pt}
\hrule height 0.6pt
\vspace{8pt}

\begin{center}
    {\fontsize{13}{15}\selectfont\textbf{BẢNG ĐÁP ÁN NHANH}}
\end{center}

\subsubsection*{1. BẢNG ĐÁP ÁN PHẦN I -- TRẮC NGHIỆM NHIỀU LỰA CHỌN (20 CÂU)}

\begin{center}
\renewcommand{\arraystretch}{1.22}
\begin{tabular}{|c|c||c|c||c|c||c|c||c|c|}
\hline
\textbf{Câu} & \textbf{Đáp án} & \textbf{Câu} & \textbf{Đáp án} & \textbf{Câu} & \textbf{Đáp án} & \textbf{Câu} & \textbf{Đáp án} & \textbf{Câu} & \textbf{Đáp án} \\ \hline
\textbf{1} & \textbf{A} & \textbf{5} & \textbf{A} & \textbf{9} & \textbf{A} & \textbf{13} & \textbf{A} & \textbf{17} & \textbf{A} \\ \hline
\textbf{2} & \textbf{B} & \textbf{6} & \textbf{B} & \textbf{10} & \textbf{B} & \textbf{14} & \textbf{B} & \textbf{18} & \textbf{B} \\ \hline
\textbf{3} & \textbf{C} & \textbf{7} & \textbf{C} & \textbf{11} & \textbf{C} & \textbf{15} & \textbf{C} & \textbf{19} & \textbf{C} \\ \hline
\textbf{4} & \textbf{D} & \textbf{8} & \textbf{D} & \textbf{12} & \textbf{D} & \textbf{16} & \textbf{D} & \textbf{20} & \textbf{D} \\ \hline
\end{tabular}
\end{center}

\subsubsection*{2. BẢNG ĐÁP ÁN PHẦN II -- CÂU HỎI ĐÚNG SAI (5 CÂU)}

\begin{center}
\renewcommand{\arraystretch}{1.22}
\begin{tabular}{|c|c|c|c|c|}
\hline
\textbf{Câu hỏi} & \textbf{Lệnh a)} & \textbf{Lệnh b)} & \textbf{Lệnh c)} & \textbf{Lệnh d)} \\ \hline
\textbf{Câu 1} & \textbf{Đúng} & \textbf{Đúng} & \textbf{Đúng} & \textbf{Sai} \\ \hline
\textbf{Câu 2} & \textbf{Đúng} & \textbf{Sai} & \textbf{Đúng} & \textbf{Đúng} \\ \hline
\textbf{Câu 3} & \textbf{Đúng} & \textbf{Đúng} & \textbf{Đúng} & \textbf{Sai} \\ \hline
\textbf{Câu 4} & \textbf{Đúng} & \textbf{Sai} & \textbf{Đúng} & \textbf{Sai} \\ \hline
\textbf{Câu 5} & \textbf{Đúng} & \textbf{Đúng} & \textbf{Đúng} & \textbf{Sai} \\ \hline
\end{tabular}
\end{center}

\subsubsection*{3. BẢNG ĐÁP ÁN PHẦN III -- CÂU TRẢ LỜI NGẮN (25 CÂU)}

\begin{center}
\renewcommand{\arraystretch}{1.22}
\begin{tabular}{|c|c||c|c||c|c||c|c||c|c|}
\hline
\textbf{Câu} & \textbf{Đáp số} & \textbf{Câu} & \textbf{Đáp số} & \textbf{Câu} & \textbf{Đáp số} & \textbf{Câu} & \textbf{Đáp số} & \textbf{Câu} & \textbf{Đáp số} \\ \hline
\textbf{1} & $-11$ & \textbf{6} & $4$ & \textbf{11} & $6$ & \textbf{16} & $1$ & \textbf{21} & $12$ \\ \hline
\textbf{2} & $1$ & \textbf{7} & $4$ & \textbf{12} & $3$ & \textbf{17} & $45$ & \textbf{22} & $3$ \\ \hline
\textbf{3} & $0$ & \textbf{8} & $6$ & \textbf{13} & $45$ & \textbf{18} & $5$ & \textbf{23} & $2$ \\ \hline
\textbf{4} & $3$ & \textbf{9} & $6$ & \textbf{14} & $0$ & \textbf{19} & $36$ & \textbf{24} & $3$ \\ \hline
\textbf{5} & $12$ & \textbf{10} & $0$ & \textbf{15} & $3$ & \textbf{20} & $6$ & \textbf{25} & $15$ \\ \hline
\end{tabular}
\end{center}

\vspace{8pt}
\hrule height 0.6pt
\vspace{8pt}

\begin{center}
    {\fontsize{13}{15}\selectfont\textbf{HƯỚNG DẪN GIẢI CHI TIẾT}}
\end{center}

\subsection*{PHẦN I. TRẮC NGHIỆM NHIỀU LỰA CHỌN (CÂU 1 -- 20)}

\begin{enumerate}
\item \textbf{Đáp án A.}  
Thay tọa độ $A(1; 0; 1)$ vào phương trình $(P)$: $3(1) - 2(0) + 1 - 4 = 0$ (thỏa mãn). Vậy $A \in (P)$.

\item \textbf{Đáp án B.}  
Mặt phẳng song song với $(\alpha): 2x - y + 5z - 7 = 0$ có dạng $2x - y + 5z + D = 0$ ($D \ne -7$).  
Vì mặt phẳng đi qua gốc tọa độ $O(0; 0; 0)$ nên $D = 0$. Phương trình là $2x - y + 5z = 0$.

\item \textbf{Đáp án C.}  
Khoảng cách từ điểm $A(2; -1; 3)$ đến mặt phẳng $(P)$:
\[
d(A, (P)) = \dfrac{|2 + 2(-1) - 2(3) + 12|}{\sqrt{1^2 + 2^2 + (-2)^2}} = \dfrac{|6|}{3} = 2.
\]

\item \textbf{Đáp án D.}  
Phương trình mặt phẳng theo đoạn chắn đi qua $M(3; 0; 0), N(0; 2; 0), P(0; 0; -1)$ là:
\[
\dfrac{x}{3} + \dfrac{y}{2} + \dfrac{z}{-1} = 1 \iff 2x + 3y - 6z - 6 = 0.
\]

\item \textbf{Đáp án A.}  
Xét tỉ lệ các hệ số: $\dfrac{2}{4} = \dfrac{-3}{-6} = \dfrac{1}{2} = \dfrac{-1}{-2} = \dfrac{1}{2}$.  
Do đó hai mặt phẳng $(P)$ và $(Q)$ trùng nhau.

\item \textbf{Đáp án B.}  
Vectơ chỉ phương của đường thẳng $AB$ là $\vec{AB} = (2 - 1; 4 - 2; 1 - (-1)) = (1; 2; 2)$.

\item \textbf{Đáp án C.}  
Đường thẳng đi qua $M(2; -3; 1)$ và có VTCP $\vec{u} = (-1; 2; 3)$ có phương trình tham số là:
\[
\begin{cases} x = 2 - t \\ y = -3 + 2t \\ z = 1 + 3t \end{cases}\quad (t \in \mathbb{R}).
\]

\item \textbf{Đáp án D.}  
Đường thẳng $d$ có VTCP $\vec{u} = (2; -2; 1)$ và $(P)$ có VTPT $\vec{n} = (2; -2; 1)$.  
Vì $\vec{u} = \vec{n}$ nên đường thẳng $d$ vuông góc với mặt phẳng $(P)$.

\item \textbf{Đáp án A.}  
Hình chiếu vuông góc của điểm $A(x_0; y_0; z_0)$ lên trục tọa độ $Oz$ có tọa độ là $(0; 0; z_0)$. Do đó $A'(0; 0; 5)$.

\item \textbf{Đáp án B.}  
Thay tọa độ $N(2; 1; 3)$ vào phương trình chính tắc của $d$:  
Ta có $\dfrac{2 - 1}{2} = \dfrac{1}{2} \ne \dfrac{1 + 1}{1} = 2$. Do đó $N \notin d$.

\item \textbf{Đáp án C.}  
VTCP của $d_1$ là $\vec{u}_1 = (1; 1; 0)$, của $d_2$ là $\vec{u}_2 = (0; 1; 1)$.
\[
\cos(d_1, d_2) = \dfrac{|\vec{u}_1 \cdot \vec{u}_2|}{|\vec{u}_1||\vec{u}_2|} = \dfrac{|1\cdot 0 + 1\cdot 1 + 0\cdot 1|}{\sqrt{1^2 + 1^2 + 0^2}\sqrt{0^2 + 1^2 + 1^2}} = \dfrac{1}{2} \implies (d_1, d_2) = 60^\circ.
\]

\item \textbf{Đáp án D.}  
Mặt phẳng $(Oxy)$ có VTPT $\vec{k} = (0; 0; 1)$ và $(Oyz)$ có VTPT $\vec{i} = (1; 0; 0)$.  
Vì $\vec{k} \cdot \vec{i} = 0$ nên góc giữa hai mặt phẳng tọa độ bằng $90^\circ$.

\item \textbf{Đáp án A.}  
VTCP của $d$ là $\vec{u} = (1; \sqrt{3}; 0)$, VTCP của trục $Ox$ là $\vec{i} = (1; 0; 0)$.
\[
\cos(d, Ox) = \dfrac{|1(1) + \sqrt{3}(0) + 0(0)|}{\sqrt{1 + 3}\cdot 1} = \dfrac{1}{2} \implies (d, Ox) = 60^\circ.
\]

\item \textbf{Đáp án B.}  
VTCP của $d$ là $\vec{u} = (2; 1; 2)$, VTPT của $(P)$ là $\vec{n} = (1; 1; 1)$.
\[
\sin(d, (P)) = \dfrac{|\vec{u} \cdot \vec{n}|}{|\vec{u}||\vec{n}|} = \dfrac{|2(1) + 1(1) + 2(1)|}{\sqrt{4 + 1 + 4}\sqrt{1 + 1 + 1}} = \dfrac{5}{3\sqrt{3}} = \dfrac{5\sqrt{3}}{9}.
\]

\item \textbf{Đáp án C.}  
$d_1 \perp d_2 \iff \vec{u}_1 \cdot \vec{u}_2 = 0 \iff 1(2) + m(1) + 2(-1) = 0 \iff m = 0$.

\item \textbf{Đáp án D.}  
Xét phương trình $x^2 + y^2 + z^2 - 2x + 4y - 6z + 1 = 0$:  
Ta có $a = 1, b = -2, c = 3, d = 1$.  
Biểu thức $a^2 + b^2 + c^2 - d = 1 + 4 + 9 - 1 = 13 > 0$. Do đó đây là phương trình của một mặt cầu.

\item \textbf{Đáp án A.}  
Bán kính mặt cầu là $R = OI = \sqrt{1^2 + (-3)^2 + 2^2} = \sqrt{14}$.  
Phương trình mặt cầu tâm $I(1; -3; 2)$ bán kính $R = \sqrt{14}$ là:
\[
(x - 1)^2 + (y + 3)^2 + (z - 2)^2 = 14.
\]

\item \textbf{Đáp án B.}  
Mặt cầu $(S)$ có tâm $I(-1; 2; -3)$. Mặt phẳng tiếp xúc với $(S)$ tại $A(2; 2; 1)$ nhận vectơ $\vec{IA} = (2 - (-1); 2 - 2; 1 - (-3)) = (3; 0; 4)$ làm vectơ pháp tuyến.

\item \textbf{Đáp án C.}  
Phương trình dạng $x^2 + y^2 + z^2 - 2ax - 2by - 2cz + d = 0 \implies -2a = -4, -2b = 6, -2c = -2 \implies a = 2, b = -3, c = 1$.  
Vậy tâm của mặt cầu là $I(2; -3; 1)$.

\item \textbf{Đáp án D.}  
Khoảng cách thẳng từ $O$ đến $M$ là độ dài đoạn $OM$:
\[
OM = \sqrt{3^2 + 4^2 + 5^2} = \sqrt{50} = 5\sqrt{2}\text{ km}.
\]
\end{enumerate}

\subsection*{PHẦN II. TRẮC NGHIỆM ĐÚNG SAI (CÂU 1 -- 5)}

\begin{enumerate}
\item \textbf{Câu 1:}
\begin{enumerate}[label=\alph*)]
    \item \textbf{Đúng.} Vì $ABCD$ là hình chữ nhật ở mặt phẳng đáy $(Oxy)$ nên $C = B + D - A = (4; 6; 0)$. Đỉnh $C'$ ngay phía trên $C$ nên $C'(4; 6; 3)$.
    \item \textbf{Đúng.} Điểm $B(4; 0; 0), D(0; 6; 0), A'(0; 0; 3)$ thuộc các trục tọa độ. Phương trình theo đoạn chắn là $\dfrac{x}{4} + \dfrac{y}{6} + \dfrac{z}{3} = 1 \iff 3x + 2y + 4z - 12 = 0$.
    \item \textbf{Đúng.} Thể tích khối hộp chữ nhật $V = AB \cdot AD \cdot AA' = 4 \times 6 \times 3 = 72$.
    \item \textbf{Sai.} Ta có $d(A, (A'BD)) = \dfrac{|-12|}{\sqrt{3^2 + 2^2 + 4^2}} = \dfrac{12}{\sqrt{29}}$\\
    và $d(C, (A'BD)) = \dfrac{|3(4) + 2(6) + 4(0) - 12|}{\sqrt{29}} = \dfrac{12}{\sqrt{29}}$. Khoảng cách từ $C$ bằng khoảng cách từ $A$, không phải gấp đôi.
\end{enumerate}

\item \textbf{Câu 2:}
\begin{enumerate}[label=\alph*)]
    \item \textbf{Đúng.} Thay $t = 0$ ta được $M(1; -2; 3) \in d$.
    \item \textbf{Sai.} Đường thẳng $d$ có VTCP $\vec{u} = (2; 1; -1)$, mặt phẳng $(P)$ có VTPT $\vec{n} = (1; 1; 3)$.  
    Ta có $\vec{u} \cdot \vec{n} = 2(1) + 1(1) + (-1)(3) = 0$. Thay điểm $M(1; -2; 3) \in d$ vào $(P)$: $1 + (-2) + 3(3) - 8 = 0$ (thỏa mãn). Vậy $d$ nằm trong $(P)$, không phải song song.
    \item \textbf{Đúng.} Vì $\vec{u} \perp \vec{n}$ và $M(1; -2; 3) \in (P)$ nên $d \subset (P)$.
    \item \textbf{Đúng.} $(Q) \perp d$ nên nhận $\vec{u} = (2; 1; -1)$ làm VTPT:  
    $2(x - 0) + 1(y - 1) - 1(z - 2) = 0 \iff 2x + y - z + 1 = 0$.
\end{enumerate}

\item \textbf{Câu 3:}
\begin{enumerate}[label=\alph*)]
    \item \textbf{Đúng.} Tâm mặt cầu là trung điểm của đoạn $AB$: $I\left(\dfrac{1+3}{2}; \dfrac{0+2}{2}; \dfrac{2+0}{2}\right) = (2; 1; 1)$.
    \item \textbf{Đúng.} $AB = \sqrt{(3-1)^2 + (2-0)^2 + (0-2)^2} = \sqrt{4 + 4 + 4} = 2\sqrt{3} \implies R = \sqrt{3}$.
    \item \textbf{Đúng.} Khoảng cách $d(I, (\alpha)) = \dfrac{|2 + 2(1) - 2(1) + 6|}{\sqrt{1^2 + 2^2 + (-2)^2}} = \dfrac{8}{3}$.
    \item \textbf{Sai.} Vì $d(I, (\alpha)) = \dfrac{8}{3} \approx 2{,}67 > \sqrt{3} \approx 1{,}732 = R$ nên mặt phẳng $(\alpha)$ không cắt mặt cầu $(S)$.
\end{enumerate}

\item \textbf{Câu 4:}
\begin{enumerate}[label=\alph*)]
    \item \textbf{Đúng.} VTCP của $d_1$ là $\vec{u}_1 = (1; -2; 2)$, của $d_2$ là $\vec{u}_2 = (2; 1; -1)$.
    \item \textbf{Sai.} $\vec{u}_1 \cdot \vec{u}_2 = 1(2) - 2(1) + 2(-1) = -2 \ne 0 \implies d_1$ không vuông góc với $d_2$.
    \item \textbf{Đúng.} $\cos(d_1, d_2) = \dfrac{|\vec{u}_1 \cdot \vec{u}_2|}{|\vec{u}_1||\vec{u}_2|} = \dfrac{|-2|}{\sqrt{9}\sqrt{6}} = \dfrac{2}{3\sqrt{6}} = \dfrac{\sqrt{6}}{9}$.
    \item \textbf{Sai.} Lấy $M_1(1; -1; 2) \in d_1, M_2(-2; 3; 0) \in d_2 \implies \vec{M_1M_2} = (-3; 4; -2)$.  
    Ta có $[\vec{u}_1, \vec{u}_2] = (0; 5; 5) \implies [\vec{u}_1, \vec{u}_2] \cdot \vec{M_1M_2} = 0(-3) + 5(4) + 5(-2) = 10 \ne 0$. Do đó $d_1$ và $d_2$ chéo nhau, không cắt nhau.
\end{enumerate}

\item \textbf{Câu 5:}
\begin{enumerate}[label=\alph*)]
    \item \textbf{Đúng.} Khoảng cách từ tâm $O$ đến $A$ là $OA = \sqrt{0^2 + 3{,}84^2 + 5{,}12^2} = \sqrt{40{,}96} = 6{,}4$ đơn vị $= 6400\text{ km}$, đúng bằng bán kính Trái Đất.
    \item \textbf{Đúng.} Khoảng cách $OS = 42$ đơn vị $= 42 \times 1000 = 42000\text{ km}$.
    \item \textbf{Đúng.} $AS = \sqrt{0^2 + (-3{,}84)^2 + (42 - 5{,}12)^2} = \sqrt{14{,}7456 + 1360{,}1344} = \sqrt{1374{,}88} \approx 37{,}08$ đơn vị $= 37080\text{ km} < 38000\text{ km}$.
    \item \textbf{Sai.} Mặt phẳng tiếp diện với Trái Đất tại $A$ có VTPT là $\vec{n} = \vec{OA} = (0; 3{,}84; 5{,}12)$.  
    Vectơ $\vec{AS} = (0; -3{,}84; 36{,}88)$.  
    Sin góc nâng $\sin\alpha = \dfrac{|\vec{n} \cdot \vec{AS}|}{|\vec{n}||\vec{AS}|} = \dfrac{174{,}08}{6{,}4 \times 37{,}08} \approx 0{,}7335 \implies \alpha \approx 47{,}2^\circ < 80^\circ$.
\end{enumerate}
\end{enumerate}

\subsection*{PHẦN III. TRẢ LỜI NGẮN (CÂU 1 -- 25)}

\begin{enumerate}
\item \textbf{Đáp số: $-11$.}  
Phương trình $(ABC)$ theo đoạn chắn: $\dfrac{x}{1} + \dfrac{y}{2} + \dfrac{z}{3} = 1 \iff 6x + 3y + 2z - 6 = 0$.  
Mặt phẳng $(P) \parallel (ABC)$ có dạng $6x + 3y + 2z + D = 0$.  
Đi qua $M(1; 1; 1) \implies 6(1) + 3(1) + 2(1) + D = 0 \iff 11 + D = 0 \iff D = -11$.

\item \textbf{Đáp số: $1$.}  
\[
d(A, (P)) = \dfrac{|3(1) - 4(3) + 14|}{\sqrt{3^2 + 0^2 + (-4)^2}} = \dfrac{|5|}{5} = 1.
\]

\item \textbf{Đáp số: $0$.}  
$\vec{AB} = (-1; 1; 1)$, $\vec{n}_Q = (1; 1; -1) \implies [\vec{AB}, \vec{n}_Q] = (-2; 0; -2) = -2(1; 0; 1)$.  
Phương trình $(P)$ nhận $\vec{n}_P = (1; 0; 1)$ làm VTPT và đi qua $A(1; 1; 0)$:
\[
1(x - 1) + 0(y - 1) + 1(z - 0) = 0 \iff x + z - 1 = 0.
\]
Dạng $Ax + By + z + D = 0 \implies A = 1, B = 0, D = -1 \implies S = 1 + 0 - 1 = 0$.

\item \textbf{Đáp số: $3$.}  
Lấy điểm $M(1; 0; 0) \in (P)$. Khoảng cách giữa hai mặt phẳng song song là:
\[
d((P), (Q)) = d(M, (Q)) = \dfrac{|1 - 2(0) + 2(0) + 8|}{\sqrt{1^2 + (-2)^2 + 2^2}} = \dfrac{9}{3} = 3.
\]

\item \textbf{Đáp số: $12$.}  
\[
d(O, (P)) = \dfrac{|m|}{\sqrt{2^2 + (-1)^2 + 2^2}} = \dfrac{m}{3} = 4 \implies m = 12.
\]

\item \textbf{Đáp số: $4$.}  
Khoảng cách từ quạt trần $I(1; 2; 4)$ đến sàn nhà nằm trên mặt phẳng tọa độ $(Oxy): z = 0$ chính là cao độ $z_I = 4\text{ mét}$.

\item \textbf{Đáp số: $4$.}  
Phương trình tham số của $d$: $x = 3 + t, y = -1 - 2t, z = 2 + t$.  
Thay vào $(P)$: $(3 + t) + 2(-1 - 2t) + (2 + t) + 1 = 0 \iff -2t + 4 = 0 \iff t = 2$.  
Tọa độ giao điểm $M(5; -5; 4) \implies T = 5 + (-5) + 4 = 4$.

\item \textbf{Đáp số: $6$.}  
Đường thẳng $d$ đi qua $B(1; 2; 3)$ và có VTCP $\vec{u} = (1; 1; 1)$. Ta có $\vec{BA} = (1; -1; 0)$.  
Vì $\vec{BA} \cdot \vec{u} = 1(1) - 1(1) + 0(1) = 0$ nên $BA \perp d$.  
Do đó hình chiếu vuông góc của $A$ lên $d$ chính là điểm $B(1; 2; 3) \implies S = 1 + 2 + 3 = 6$.

\item \textbf{Đáp số: $6$.}  
Điểm đối xứng của $A(2; 1; -3)$ qua mặt phẳng $(Oxy)$ là $A'(2; 1; 3)$.  
Độ dài $AA' = |-3 - 3| = 6$.

\item \textbf{Đáp số: $0$.}  
$d_1$ có VTCP $\vec{u}_1 = (2; 3; 4)$, $d_2$ có VTCP $\vec{u}_2 = (4; 6; 8) = 2\vec{u}_1$.  
Điểm $M_1(1; 2; 3) \in d_1$. Thay vào phương trình $d_2$:  
$\dfrac{1 - 3}{4} = \dfrac{2 - 5}{6} = \dfrac{3 - 7}{8} = -\dfrac{1}{2}$ (thỏa mãn), suy ra $M_1 \in d_2$.  
Vậy $d_1$ trùng với $d_2$. Theo quy ước đề bài ghi đáp số bằng $0$.

\item \textbf{Đáp số: $6$.}  
Tọa độ trọng tâm $G$ của tam giác $ABC$:
\[
x_G = \dfrac{1 + 2 + 3}{3} = 2;\quad y_G = \dfrac{2 + 1 + 3}{3} = 2;\quad z_G = \dfrac{3 + 0 + 3}{3} = 2.
\]
Vậy $x_G + y_G + z_G = 2 + 2 + 2 = 6$.

\item \textbf{Đáp số: $3$.}  
$d_1$ đi qua $A(1; 0; 2)$ có VTCP $\vec{u}_1 = (1; 0; 0)$.  
$d_2$ đi qua $B(0; 0; -1)$ có VTCP $\vec{u}_2 = (0; 1; 0)$.  
Ta có $[\vec{u}_1, \vec{u}_2] = (0; 0; 1)$ và $\vec{AB} = (-1; 0; -3)$.  
Khoảng cách ngắn nhất là $d(d_1, d_2) = \dfrac{|[\vec{u}_1, \vec{u}_2] \cdot \vec{AB}|}{|[\vec{u}_1, \vec{u}_2]|} = \dfrac{|-3|}{1} = 3$.

\item \textbf{Đáp số: $45$.}  
$\vec{u} = (1; 1; 0)$, VTCP của trục $Oy$ là $\vec{j} = (0; 1; 0)$.
\[
\cos(d, Oy) = \dfrac{|1\cdot 0 + 1\cdot 1 + 0\cdot 0|}{\sqrt{1 + 1}\cdot 1} = \dfrac{\sqrt{2}}{2} \implies (d, Oy) = 45^\circ.
\]

\item \textbf{Đáp số: $0$.}  
$\vec{u} = (\sqrt{3}; 1; 0)$, VTPT của $(P)$ là $\vec{k} = (0; 0; 1)$.  
Ta có $\vec{u} \cdot \vec{k} = 0 \implies \sin(d, (P)) = 0 \implies$ góc giữa đường thẳng và mặt phẳng bằng $0^\circ$ ($d \parallel (P)$).

\item \textbf{Đáp số: $3$.}  
$\vec{n}_P = (1; 1; 0) \implies |\vec{n}_P| = \sqrt{2}$; $\vec{n}_Q = (0; 1; 1) \implies |\vec{n}_Q| = \sqrt{2}$.
\[
\cos((P), (Q)) = \dfrac{|1\cdot 0 + 1\cdot 1 + 0\cdot 1|}{\sqrt{2}\sqrt{2}} = \dfrac{1}{2} \implies a = 1, b = 2 \implies a + b = 3.
\]

\item \textbf{Đáp số: $1$.}  
$\vec{u}_1 = (1; 0; m) \implies |\vec{u}_1| = \sqrt{1 + m^2}$; $\vec{u}_2 = (1; 1; 0) \implies |\vec{u}_2| = \sqrt{2}$.
\[
\cos 60^\circ = \dfrac{1}{\sqrt{2(1 + m^2)}} = \dfrac{1}{2} \iff 2(1 + m^2) = 4 \iff m^2 = 1 \implies m = 1.
\]

\item \textbf{Đáp số: $45$.}  
VTPT của tấm pin là $\vec{n} = (1; 1; \sqrt{2})$, VTPT của mặt phẳng sàn là $\vec{k} = (0; 0; 1)$.
\[
\cos\alpha = \dfrac{|\sqrt{2}|}{\sqrt{1 + 1 + 2}\cdot 1} = \dfrac{\sqrt{2}}{2} \implies \alpha = 45^\circ.
\]

\item \textbf{Đáp số: $5$.}  
$a = 2, b = -1, c = 0, d = -20 \implies R = \sqrt{2^2 + (-1)^2 + 0^2 - (-20)} = \sqrt{25} = 5$.

\item \textbf{Đáp số: $36$.}  
Bán kính mặt cầu là khoảng cách từ tâm $I$ đến tiếp diện $(P)$:
\[
R = d(I, (P)) = \dfrac{|2(2) - 2(1) + (-1) + 8|}{\sqrt{2^2 + (-2)^2 + 1^2}} = \dfrac{9}{3} = 3.
\]
Diện tích toàn phần chia cho $\pi$: $\dfrac{S}{\pi} = 4R^2 = 4(3^2) = 36$.

\item \textbf{Đáp số: $6$.}  
Mặt cầu có tâm $O(0; 0; 0)$ và bán kính $R = \sqrt{13}$.  
Khoảng cách $d = d(O, (P)) = |-2| = 2$.  
Bán kính đường tròn giao tuyến: $r = \sqrt{R^2 - d^2} = \sqrt{13 - 4} = 3$.  
Chu vi đường tròn chia cho $\pi$: $\dfrac{C}{\pi} = 2r = 6$.

\item \textbf{Đáp số: $12$.}  
Điểm trên đường ray có tọa độ $M(2; -5 + t; 2)$.  
Khoảng cách từ $I(2; 3; 2)$ đến $M$ là $IM = |(-5 + t) - 3| = |t - 8|$.  
Trong vùng phủ sóng: $IM \le 6 \iff |t - 8| \le 6 \iff 2 \le t \le 14$.  
Quãng đường robot di chuyển trong vùng sóng là: $14 - 2 = 12\text{ mét}$.

\item \textbf{Đáp số: $3$.}  
Mặt cầu $(S)$ có tâm $I(1; -1; 0)$ và bán kính $R = 2$.  
Mặt phẳng $(P)$ tiếp xúc với $(S)$ khi:
\[
d(I, (P)) = R \iff \dfrac{|2(1) - (-1) + 2(0) + m|}{\sqrt{2^2 + (-1)^2 + 2^2}} = 2 \iff \dfrac{|m + 3|}{3} = 2 \iff |m + 3| = 6.
\]
Vì $m > 0$ nên $m + 3 = 6 \implies m = 3$.

\item \textbf{Đáp số: $2$.}  
Điểm $M \in Ox \implies M(x; 0; 0)$.  
Ta có $MA^2 = (x - 1)^2 + 1 + 1 = x^2 - 2x + 3$.  
$MB^2 = (x - 3)^2 + 9 + 9 = x^2 - 6x + 27$.  
$f(x) = MA^2 + MB^2 = 2x^2 - 8x + 30$.  
Tam thức bậc hai đạt giá trị nhỏ nhất tại hoành độ đỉnh $x = -\dfrac{-8}{2(2)} = 2$.

\item \textbf{Đáp số: $3$.}  
Đường thẳng đi qua $A(3; 2; 2)$ vuông góc với $(P)$ có phương trình: $x = 3 + t, y = 2 + t, z = 2 + t$.  
Thay vào phương trình $(P)$: $(3 + t) + (2 + t) + (2 + t) - 1 = 0 \iff 3t + 6 = 0 \iff t = -2$.  
Tọa độ hình chiếu là $H(1; 0; 0)$.  
Tọa độ điểm đối xứng: $A' = 2H - A = (2(1) - 3; 2(0) - 2; 2(0) - 2) = (-1; -2; -2)$.  
Khoảng cách $OA' = \sqrt{(-1)^2 + (-2)^2 + (-2)^2} = \sqrt{1 + 4 + 4} = 3$.

\item \textbf{Đáp số: $15$.}  
Mặt cầu $(S)$ có tâm $I(0; 0; 10)$ và bán kính $R = 5$.  
Điểm $M(x; y; z) \in (S)$ thỏa mãn $(z - 10)^2 \le 25 \iff -5 \le z - 10 \le 5 \iff 5 \le z \le 15$.  
Do đó độ cao lớn nhất mà đầu cần cẩu có thể vươn tới là $z_{\max} = 15\text{ mét}$.
\end{enumerate}

\end{document}
